Modern systems are becoming increasingly heterogeneous in hardware.

Yet programming on heterogeneous systems remains daunting for software developers, who often lack the means to reason about how their design choices will affect performance on diverse architectures.

Prior works attemp to tackle this challenge in two separate landscapes: static compile-time reasoning and dynamic runtime prediction.

% Prior work has largely  this challenge in two separate landscapes: static compile-time reasoning and dynamic runtime prediction. Static approaches help developers reason about performance before execution, such as compiler cost models [cite], analytical models [cite], and autotuning [cite]. Runtime approaches predict behavior under real workloads, such as simulation [cite], profiling [cite], and machine learning-based modeling [cite]. Some techniques bridge the two, notably profile-guided optimization [cite] and JIT compilation [cite], which feed runtime observations back into the compiler. However, these approaches optimize code transparently rather than helping developers reason about the performance implications of their design choices on heterogeneous hardware.

In this paper, we present \sysname, a system that unifies these two perspectives by providing a framework for performance prediction on heterogeneous systems.